\documentclass[a4paper,
11pt,  reqno]{amsart}
\usepackage[margin=1.2in,marginparwidth=1.5cm, marginparsep=0.5cm]{geometry}
\usepackage{graphicx} % Required for inserting images
\usepackage{amsmath, amssymb, amsthm}
\usepackage{mathtools}

\title{Exercícios de ``Introdução às Equações Diferenciais''  IMPA-TECH 2026-01}

\begin{document}

\maketitle

\section*{Simulado 1}

\vspace{5mm}

\begin{enumerate}

\item Resolva os seguintes problemas de valor inicial. Em cada item, indique explicitamente o maior intervalo em que a solução está definida.

\begin{enumerate}
    \item 
    \[
        y' + 2y = e^{-t}, \qquad y(0)=1.
    \]

    \item 
    \[
        y' = t(1+y^2), \qquad y(0)=0.
    \]

    \item 
    \[
        y' = \frac{t+y}{t}, \qquad y(1)=0, \qquad t>0.
    \]
    \emph{Dica:} tente reconhecer uma expressão simples cuja derivada aparece na equação.
\end{enumerate}

\vspace{2mm}

\item Considere o problema de valor inicial
\[
    \frac{dy}{dt}=3y^{2/3},
    \qquad
    y(t_0)=y_0.
\]

\begin{enumerate}
    \item No ponto $(t_0,y_0)=(0,1)$, determine se existe solução e se ela é única. Justifique.

    \item Faça o mesmo para o ponto $(t_0,y_0)=(0,0)$. Caso a solução não seja única, exiba duas soluções distintas.

    \item Explique por que os dois casos acima têm comportamentos diferentes à luz do teorema de existência e unicidade.
\end{enumerate}


\vspace{2mm}

\item Considere a equação diferencial
\[
    t^2y''-2ty'+2y=t^3,
    \qquad
    t>0.
\]

\begin{enumerate}
    \item Verifique que $y_1(t)=t$ é uma solução da equação homogênea associada.

    \item Encontre uma segunda solução $y_2$ da equação homogênea que seja linearmente independente de $y_1$.

    \item Encontre a solução geral da equação não homogênea.

\end{enumerate}

\vspace{2mm}


\item Seja $a:[0,1]\to(0,+\infty)$ uma função contínua, e seja $\lambda\in\mathbb{R}$. Considere a equação
\[
    x''+\lambda a(t)x=0.
\]
Suponha que exista uma solução não identicamente nula $x$ tal que
\[
    x(0)=0,
    \qquad
    x'(1)=0.
\]

\begin{enumerate}
    \item Mostre que, se $\lambda\leq 0$, então tal solução não pode existir.

    \item Exiba um exemplo com $\lambda>0$ para o qual existe uma solução não identicamente nula satisfazendo
    \[
        x(0)=0,
        \qquad
        x'(1)=0.
    \]

\end{enumerate}

\vspace{6mm}
\newpage

\item Considere os itens abaixo:

\begin{enumerate}
    \item Mostre que não existe uma equação linear homogênea de segunda ordem com coeficientes contínuos em todo $\mathbb{R}$,
\[
    y''+p(t)y'+q(t)y=0,
\]
que possua simultaneamente as duas funções
\[
    y_1(t)=\sin(t^2)
    \qquad\text{e}\qquad
    y_2(t)=\cos(t^2)
\]
como soluções em todo $\mathbb{R}$.

    \item O resultado do item anterior vale ainda se substituirmos $\mathbb{R}$ por $(0,+\infty)$? 
\end{enumerate}




\end{enumerate}

\end{document}